\chapter{HIV Screening}

\section*{What Is This Test?}
HIV screening tests detect human immunodeficiency virus (HIV) infection through identification of HIV-1 and HIV-2 antibodies, HIV-1 p24 antigen, or HIV nucleic acid. HIV is a retrovirus (genus \textit{Lentivirus}) that infects CD4+ T lymphocytes, macrophages, and dendritic cells via the CD4 receptor and CCR5/CXCR4 co-receptors. Untreated HIV infection progresses through stages: acute (seroconversion) illness (fever, lymphadenopathy, pharyngitis, rash, myalgia 2--4 weeks post-infection), clinical latency (asymptomatic, 2--10+ years with progressive CD4 decline), and AIDS (CD4 <200 cells/$\mu$L or an AIDS-defining condition). HIV screening in the United States uses a 4th-generation antigen/antibody combination immunoassay (detecting HIV-1 p24 antigen and HIV-1/2 antibodies simultaneously), which reduces the window period to approximately 15--20 days from infection. Reactive screening tests are confirmed with an HIV-1/HIV-2 antibody differentiation immunoassay, and if that is indeterminate or negative, HIV-1 nucleic acid testing (NAT) resolves the status. This CDC-recommended algorithm ensures high sensitivity and specificity. Rapid point-of-care tests (OraQuick, INSTI) detect antibodies only and have sensitivity comparable to laboratory-based 4th-generation tests but with a longer window period (~23--30 days).

\section*{Alternative Names}
HIV Test; HIV 4th Generation Test; HIV Antigen/Antibody Test; HIV-1/2 Antibody Test; p24 Antigen Test; HIV ELISA; HIV Rapid Test; HIV Screening; HIV Serology; AIDS Test; OraQuick

\section*{Principle}
4th-generation HIV-1/2 antigen/antibody combination immunoassay: A chemiluminescent microparticle immunoassay (CMIA) or enzyme immunoassay (EIA) simultaneously detects HIV-1 p24 antigen (capsid protein released during acute infection) and IgG/IgM antibodies against HIV-1 (groups M, O) and HIV-2. Antigen is detectable within 10--20 days of infection, antibodies within 20--30 days. Sensitivity >99.7\%, specificity >99.3\%. If the screening assay is reactive, the specimen is automatically reflexed to an HIV-1/HIV-2 antibody differentiation immunoassay (e.g., Geenius HIV 1/2 Confirmatory Assay, Bio-Rad Multispot), which distinguishes HIV-1 from HIV-2 antibodies and serves as the first step in the confirmation algorithm. If the differentiation assay is negative or indeterminate, HIV-1 nucleic acid testing (qualitative or quantitative HIV-1 RNA PCR) is performed to determine whether acute HIV infection (p24 antigen positive, antibodies not yet detectable) is present. The 3-step algorithm (screening $\rightarrow$ differentiation $\rightarrow$ NAT) has supplanted the older Western blot algorithm, which had a longer window period and detected antibodies only \cite{branson2014laboratory}. Rapid tests: Lateral flow immunochromatographic assays detecting HIV-1/2 antibodies in oral fluid (OraQuick), fingerstick blood, or venous blood; results available in 20--40 minutes; sensitivity 99.3--99.7\%, specificity 99.8--99.9\%.

\section*{History}
HIV was first recognized in 1981 with reports of \textit{Pneumocystis} pneumonia and Kaposi sarcoma in previously healthy men who have sex with men (MSM). HIV-1 was isolated by Luc Montagnier and Françoise Barré-Sinoussi in 1983 (Nobel Prize 2008). The first HIV antibody test (ELISA) was licensed in 1985 \cite{alexander2016human}. 2nd-generation tests (recombinant antigen-based) improved sensitivity in the 1990s. 3rd-generation tests (IgM-sensitive antibody detection) reduced the window period to approximately 20--25 days. 4th-generation combination antigen/antibody assays were introduced in the late 2000s and recommended by the CDC as the screening standard in 2014. The first rapid test (OraQuick) was FDA-approved in 2002; FDA approved the first over-the-counter HIV self-test (OraQuick In-Home HIV Test) in 2012. The CDC's 2014 revised testing algorithm eliminated Western blot confirmatory testing in favor of the HIV-1/2 differentiation assay plus NAT workflow. Antiretroviral therapy (ART), introduced in 1987 (zidovudine/AZT) and revolutionized with combination ART in 1996, has transformed HIV from a fatal disease to a manageable chronic condition. The concept of U=U (Undetectable = Untransmittable), supported by HPTN 052, PARTNER, and Opposites Attract studies, has been CDC-endorsed since 2017.

\section*{Why This Test Is Ordered}
\begin{itemize}
  \item Routine HIV screening: The CDC recommends at least one-time HIV screening for all individuals aged 13--64 years in healthcare settings, with opt-out testing (the patient is informed that testing will be performed unless they decline) \cite{cdc2006revised}
  \item Annual or more frequent screening for individuals at higher risk: MSM (every 3--6 months), people who inject drugs, commercial sex workers, individuals with multiple sexual partners, sexual partners of HIV-positive individuals, and those diagnosed with another STI
  \item Evaluation of individuals presenting with symptoms suggestive of acute HIV seroconversion syndrome: fever, pharyngitis, lymphadenopathy, maculopapular rash (typically on the trunk), myalgia/arthralgia, headache, oral/genital ulcers, or aseptic meningitis, 2--4 weeks after a high-risk exposure \cite{daar2001diagnosis}
  \item Universal prenatal screening: All pregnant women should be screened for HIV at the first prenatal visit (opt-out approach); repeat testing in the third trimester is recommended for women at high risk of acquisition, in areas with high HIV incidence, or for those who declined initial testing
  \item Testing after known or suspected occupational or non-occupational exposure to HIV (needlestick, sexual assault): baseline HIV test immediately after exposure, with follow-up testing at 4--6 weeks and 3 months
  \item Pre-exposure prophylaxis (PrEP) initiation and monitoring: Documentation of negative HIV status before PrEP initiation (TDF/FTC or TAF/FTC), with quarterly HIV screening plus assessment of renal function and STI screening while on PrEP
  \item Preoperative or pre-procedural testing in some institutional protocols, particularly for organ or tissue donors
\end{itemize}

\section*{Primary vs Secondary Test}
HIV screening is a primary test. The 4th-generation antigen/antibody combination assay is the recommended primary screening test in the United States for all individuals aged 13--64 and for those at ongoing risk. A reactive screening test requires confirmation (differentiation assay plus NAT if needed). HIV-1 RNA qualitative or quantitative testing is a primary test in specific settings: screening of potential blood/organ donors, diagnosis of acute HIV infection when seroconversion is suspected but 4th-generation test may be negative or indeterminate (first 10--15 days).

\section*{How This Test Is Performed}
Venous blood is collected in an SST (serum separator), EDTA (plasma), or heparin tube. For 4th-generation testing, either serum or plasma is acceptable; approximately 5 mL of whole blood is sufficient. For rapid testing, oral fluid (OraQuick: swab of the upper and lower outer gums), fingerstick whole blood, or venous whole blood may be used depending on the specific kit. Specimens for 4th-generation testing should be centrifuged and serum/plasma separated within 4 hours; specimens can be refrigerated at 2--8\textdegree{}C for up to 7 days or frozen at -20\textdegree{}C for longer storage. The initial screening test result is typically available within 1--24 hours. If the screening is reactive, reflex testing (differentiation assay and, if needed, NAT) is performed by the laboratory on the same specimen (reflex algorithm) without requiring a second blood draw in most cases.

\section*{What Preparation Is Needed}
No special patient preparation (no fasting required). No dietary or activity restrictions. Patients should be counseled about the testing procedure and the meaning of possible results (including the window period). For occupational exposure, source patient testing requires informed consent in most jurisdictions (varies by state).

\section*{Contraindications and Precautions}
The 4th-generation antigen/antibody assay has a window period of approximately 15--20 days from HIV acquisition. A negative result within the window period does NOT exclude HIV infection. If recent high-risk exposure (within the last 4--6 weeks) is reported, the patient should be counseled about the window period and offered repeat testing at 4--6 weeks and 3 months post-exposure using a 4th-generation test, or immediate HIV-1 RNA testing if acute infection is clinically suspected. False-positive screening results occur in 0.1--0.4\% of tests; all reactive screening results must be confirmed using the CDC algorithm. A reactive screen with negative confirmatory tests (differentiation and NAT) represents a false-positive screen, and no further HIV testing is needed unless a new exposure occurs. False-negative rapid tests (oral fluid) have been reported during acute HIV infection; 4th-generation laboratory-based tests are preferred if acute infection is suspected. The INSTI rapid test (bioLytical) detects antibodies only and should not be used as a standalone screening test in high-risk populations without a confirmatory algorithm. HIV testing should be voluntary and accompanied by pre-test counseling (or at minimum, the provision of information about the test) and post-test counseling.

\section*{Risks}
The primary risk is not medical but psychological/social: stress associated with a reactive (potentially false-positive) screening result and the implications of HIV diagnosis. Confidentiality protections under HIPAA and state laws are critical. Venipuncture carries minimal physical risk. Disclosure to sexual partners is legally mandated in some jurisdictions. Pre-test and post-test counseling should address these issues.

\section*{What Happens After the Test}
4th-generation screening test results typically available in 1--24 hours. Rapid test results available in 20--40 minutes. A nonreactive result: no further testing indicated unless exposure occurred within the window period. A reactive result: the laboratory reflexively performs the differentiation assay on the same specimen. If the differentiation assay is positive for HIV-1 or HIV-2 antibodies, the diagnosis is confirmed, and the patient should be linked to HIV care for ART initiation. If the differentiation assay is negative or indeterminate, HIV-1 RNA NAT is performed. A positive NAT confirms acute HIV-1 infection (seronegative, RNA positive). A negative NAT with a reactive screen and negative/indeterminate differentiation assay indicates a false-positive screening test. Confirmed positive results must be reported to the state health department (HIV is a nationally notifiable condition).

\section*{Understanding Results}

\subsection*{Below Normal}
\begin{itemize}
  \item \textbf{Nonreactive (negative) 4th-generation HIV test:} No HIV-1 p24 antigen or HIV-1/2 antibodies detected. This result effectively rules out HIV infection acquired 4 or more weeks prior. Key caveat: A negative result within the window period (first 15--30 days after exposure) does NOT exclude acute HIV infection. If exposure occurred within the preceding 4 weeks, the test should be repeated at 4--6 weeks and at 3 months post-exposure. HIV-1 RNA testing should be performed if acute seroconversion syndrome is suspected
  \item \textbf{Nonreactive rapid test:} Same interpretation but longer window period (23--30 days). Individuals at high risk with a negative rapid test should be offered 4th-generation laboratory testing or NAT if acute infection is suspected
\end{itemize}

\subsection*{Above Normal}
\begin{itemize}
  \item \textbf{Reactive (positive) 4th-generation test followed by positive HIV-1/HIV-2 differentiation assay:} Confirms HIV infection. HIV-1 positive: the vast majority of HIV infections worldwide. HIV-2 positive: rare, primarily in West Africa; important to distinguish because HIV-2 is intrinsically resistant to non-nucleoside reverse transcriptase inhibitors (NNRTIs: efavirenz, rilpivirine). Linkage to HIV care and immediate ART initiation is recommended regardless of CD4 count (DHHS guidelines \cite{dhhs2023guidelines})
  \item \textbf{Reactive 4th-generation test with negative/indeterminate differentiation assay and positive HIV-1 RNA NAT:} Confirms acute (primary) HIV-1 infection. p24 antigen is present, but antibodies have not yet developed (window period). Patients are highly viremic (>100,000 copies/mL typical) and at greatest risk of transmitting HIV. ART should be initiated immediately
  \item \textbf{Reactive 4th-generation test with negative differentiation assay and negative HIV-1 RNA NAT:} False-positive screening test. No HIV infection. The patient should be reassured and counseled; no further HIV testing is required unless a new exposure occurs. This pattern may also represent very early HIV-2 infection (HIV-2 RNA NAT not typically performed unless there is epidemiological risk)
  \item \textbf{CD4 count low and/or HIV RNA high:} Confirmation of HIV diagnosis is followed by baseline laboratory evaluation (see Follow-up Tests)
\end{itemize}

\section*{Normal Results}
\textbf{HIV-1/2 Antigen/Antibody: Nonreactive (Negative).} No HIV infection detected. For individuals with documented nonreactive status and no recent (within 4--6 weeks) exposure, no further testing required.

\section*{Follow-up Tests}
\begin{itemize}
  \item \textbf{HIV-1/2 antibody differentiation immunoassay (Geenius, Multispot):} Reflex test performed automatically on all reactive 4th-generation screening tests; distinguishes HIV-1 from HIV-2 antibodies
  \item \textbf{HIV-1 RNA quantitative (viral load):} If acute HIV is suspected (seroconversion symptoms, recent high-risk exposure) or screening is reactive with negative differentiation assay. Baseline HIV RNA measurement is also standard of care for all newly diagnosed patients as part of the initial evaluation
  \item \textbf{CD4 T-lymphocyte count (absolute and percentage):} Staging test for all newly diagnosed HIV-infected patients. CD4 count >500 cells/$\mu$L (Stage 1), 200--499 (Stage 2), <200 (Stage 3, AIDS). CD4 monitoring guides opportunistic infection (OI) prophylaxis: CD4 <200: PCP prophylaxis (TMP-SMX); CD4 <100: toxoplasmosis prophylaxis if seropositive; CD4 <50: MAC prophylaxis (azithromycin)
  \item \textbf{HIV genotype resistance testing (integrase, protease, RT):} Recommended at entry into care for all treatment-naive patients to detect transmitted drug resistance; guides initial ART selection
  \item \textbf{HLA-B*5701 testing:} Before initiation of abacavir-containing ART; HLA-B*5701-positive patients have a ~50\% risk of abacavir hypersensitivity reaction and should not receive abacavir
  \item \textbf{Coreceptor tropism testing:} If maraviroc is being considered as part of ART; determines whether the virus uses CCR5, CXCR4, or both (dual/mixed tropism)
  \item \textbf{Baseline laboratories (CBC, comprehensive metabolic panel, lipid panel, urinalysis):} To assess renal/hepatic function, bone marrow status, and cardiovascular risk before ART initiation
\end{itemize}

\section*{References}
\bibliographystyle{unsrtnat}
\bibliography{references}

\clearpage
\section*{Regulatory Status and Authorization Date}
Regulatory status applies to the specific assay, instrument, manufacturer, and intended use rather than to the test name alone. No single approval date can be assigned to this test category. For a commercial assay, verify the current product in the relevant regulator's database: the U.S. Food and Drug Administration (FDA) clearance, approval, or authorization; India's Central Drugs Standard Control Organisation (CDSCO) licence or permission; the European Union In Vitro Diagnostic Regulation (EU IVDR) CE certificate issued through the applicable conformity-assessment route; or the United Kingdom Medicines and Healthcare products Regulatory Agency (MHRA) registration and UKCA/CE status. Laboratory-developed tests may not have an individual FDA, CDSCO, EU IVDR, or MHRA approval date.
\clearpage
